% --- Archon physics formalization source begin ---
% archon:physics
% archon:covers PhyXMiniProblems/problem_phyx_mini_0635.lean
% archon:source-report reports/phyx_mini/problem_phyx_mini_0635.source.json

\chapter{Physics problem phyx\_mini\_0635}
\label{ch:PhyXMiniProblems_problem_phyx_mini_0635}

\paragraph{Problem source.}
\#\# Physical scenario

An alpha particle is shot with a speed of \$2.0 \textbackslash{}times 10\textasciicircum{}7\textbackslash{};m/s\$ directly toward the nucleus of a gold atom.

\#\# Question

What is the distance of closest approach to the nucleus?

\#\# Answer choices

- A: A: \$2.2 \textbackslash{}times 10\textasciicircum{}\{-14\}\textbackslash{};m\$

- B: B: \$2.7 \textbackslash{}times 10\textasciicircum{}\{-14\}\textbackslash{};m\$

- C: C: \$3.1 \textbackslash{}times 10\textasciicircum{}\{-14\}\textbackslash{};m\$

- D: D: \$3.4 \textbackslash{}times 10\textasciicircum{}\{-14\}\textbackslash{};m\$

\#\# Figure caption (auxiliary; use the image as primary evidence)

The image is a diagram illustrating the interaction between an alpha particle and a gold nucleus.

1. **Before Interaction:**
   - An alpha particle, labeled with charge \textbackslash{}( q\_\textbackslash{}alpha = 2e \textbackslash{}), is shown on the left.
   - It is moving towards a gold nucleus with an initial velocity \textbackslash{}( v\_i \textbackslash{}).
   - The gold nucleus is on the right, labeled with charge \textbackslash{}( q\_\{\textbackslash{}text\{Au\}\} = 79e \textbackslash{}).

2. **After Interaction:**
   - The alpha particle is depicted close to the gold nucleus, at a distance labeled \textbackslash{}( r\_\{\textbackslash{}text\{min\}\} \textbackslash{}).
   - At this closest approach, the alpha particle's velocity \textbackslash{}( v\_f = 0 \textbackslash{}).
   - Text at the bottom states: "When the \textbackslash{}(\textbackslash{}alpha\textbackslash{}) particle is at its closest approach to the gold nucleus, its speed is zero."

The diagram appears to be illustrating the concept of closest approach and how the kinetic energy of the alpha particle is turned into potential energy as it approaches the nucleus.

\#\# Dataset metadata

Nuclear Physics; Physical Model Grounding Reasoning, Multi-Formula Reasoning

\paragraph{Recorded answer/context.}
B: B: \$2.7 \textbackslash{}times 10\textasciicircum{}\{-14\}\textbackslash{};m\$

\paragraph{Figure/image path.}
/root/proposal\_for\_physic/science-mango/phyx\_mini\_run/phyx\_data/test\_image/635.png

\paragraph{Formalization target.}
create a compiling Lean file with sorry bodies at `PhyXMiniProblems/problem\_phyx\_mini\_0635.lean`. The Lean declarations must preserve the physical quantities, dimensions or dimensional roles, figure labels, governing-law hypotheses, and final relation expressed by this problem.
Use Mathlib/Physlib names found through LeanExplore where available. If a physics API is missing, introduce faithful local abstractions rather than scalar placeholder aliases.

\begin{theorem}[Physics formalization target]
\label{thm:physics:phyx_mini_0635:target}
\lean{PhyXMiniProblems.ProblemPhyXMini0635.problem_phyx_mini_0635}
\uses{def:physics:phyx-mini-0635:phyxminiproblems-problemphyxmini0635-lengthinmeters, def:physics:phyx-mini-0635:phyxminiproblems-problemphyxmini0635-massinkilograms, def:physics:phyx-mini-0635:phyxminiproblems-problemphyxmini0635-chargeincoulombs, def:physics:phyx-mini-0635:phyxminiproblems-problemphyxmini0635-speedinmeterspersecond, def:physics:phyx-mini-0635:phyxminiproblems-problemphyxmini0635-alphagoldclosestapproachsetup, def:physics:phyx-mini-0635:phyxminiproblems-problemphyxmini0635-matchesproblemscenarioanddata, def:physics:phyx-mini-0635:phyxminiproblems-problemphyxmini0635-matchesprimaryfigure, def:physics:phyx-mini-0635:phyxminiproblems-problemphyxmini0635-usesreferencealphagoldconstants, def:physics:phyx-mini-0635:phyxminiproblems-problemphyxmini0635-hasphysicalalphagoldparameters, def:physics:phyx-mini-0635:phyxminiproblems-problemphyxmini0635-satisfiesclassicalkineticenergylaw, def:physics:phyx-mini-0635:phyxminiproblems-problemphyxmini0635-satisfiesrepulsivecoulombpotentiallaw, def:physics:phyx-mini-0635:phyxminiproblems-problemphyxmini0635-satisfiesmechanicalenergyconservation, def:physics:phyx-mini-0635:phyxminiproblems-problemphyxmini0635-recordeddatasetanswer, def:physics:phyx-mini-0635:phyxminiproblems-problemphyxmini0635-isuniquematchingdisplayedclosestapproach}
The assigned autoformalize agent should translate this physics problem into Lean declarations in the covered file, with theorem and lemma proof bodies written as `by sorry`.
\end{theorem}
\begin{proof}
This is an autoformalization task, not a proof task. Produce statements that can later be proved by the physics prover without weakening the physical model.
\end{proof}

% --- Archon declaration topology begin ---
\section{Declaration topology}
\begin{definition}[Length Quantity]
  \label{def:physics:phyx-mini-0635:phyxminiproblems-problemphyxmini0635-lengthquantity}
  \lean{PhyXMiniProblems.ProblemPhyXMini0635.LengthQuantity}
  A nonnegative, unit-independent physical length
\end{definition}
\begin{proof}
  This is the declared model definition of Length Quantity.
\end{proof}
\begin{definition}[Mass Quantity]
  \label{def:physics:phyx-mini-0635:phyxminiproblems-problemphyxmini0635-massquantity}
  \lean{PhyXMiniProblems.ProblemPhyXMini0635.MassQuantity}
  A nonnegative, unit-independent physical mass
\end{definition}
\begin{proof}
  This is the declared model definition of Mass Quantity.
\end{proof}
\begin{definition}[Charge Magnitude Quantity]
  \label{def:physics:phyx-mini-0635:phyxminiproblems-problemphyxmini0635-chargemagnitudequantity}
  \lean{PhyXMiniProblems.ProblemPhyXMini0635.ChargeMagnitudeQuantity}
  A nonnegative magnitude of electric charge
\end{definition}
\begin{proof}
  This is the declared model definition of Charge Magnitude Quantity.
\end{proof}
\begin{definition}[length Readout]
  \label{def:physics:phyx-mini-0635:phyxminiproblems-problemphyxmini0635-lengthreadout}
  \lean{PhyXMiniProblems.ProblemPhyXMini0635.lengthReadout}
  \uses{def:physics:phyx-mini-0635:phyxminiproblems-problemphyxmini0635-lengthquantity}
  Read a physical length in a selected Physlib length unit
\end{definition}
\begin{proof}
  This is the declared model definition of length Readout.
\end{proof}
\begin{definition}[length In Meters]
  \label{def:physics:phyx-mini-0635:phyxminiproblems-problemphyxmini0635-lengthinmeters}
  \lean{PhyXMiniProblems.ProblemPhyXMini0635.lengthInMeters}
  \uses{def:physics:phyx-mini-0635:phyxminiproblems-problemphyxmini0635-lengthquantity, def:physics:phyx-mini-0635:phyxminiproblems-problemphyxmini0635-lengthreadout}
  SI metre readout of a physical length
\end{definition}
\begin{proof}
  This is the declared model definition of length In Meters.
\end{proof}
\begin{definition}[length In Femtometers]
  \label{def:physics:phyx-mini-0635:phyxminiproblems-problemphyxmini0635-lengthinfemtometers}
  \lean{PhyXMiniProblems.ProblemPhyXMini0635.lengthInFemtometers}
  \uses{def:physics:phyx-mini-0635:phyxminiproblems-problemphyxmini0635-lengthquantity, def:physics:phyx-mini-0635:phyxminiproblems-problemphyxmini0635-lengthreadout}
  Femtometre readout of a nuclear-scale physical length
\end{definition}
\begin{proof}
  This is the declared model definition of length In Femtometers.
\end{proof}
\begin{definition}[mass In Kilograms]
  \label{def:physics:phyx-mini-0635:phyxminiproblems-problemphyxmini0635-massinkilograms}
  \lean{PhyXMiniProblems.ProblemPhyXMini0635.massInKilograms}
  \uses{def:physics:phyx-mini-0635:phyxminiproblems-problemphyxmini0635-massquantity}
  SI kilogram readout of a physical mass
\end{definition}
\begin{proof}
  This is the declared model definition of mass In Kilograms.
\end{proof}
\begin{definition}[charge Readout]
  \label{def:physics:phyx-mini-0635:phyxminiproblems-problemphyxmini0635-chargereadout}
  \lean{PhyXMiniProblems.ProblemPhyXMini0635.chargeReadout}
  \uses{def:physics:phyx-mini-0635:phyxminiproblems-problemphyxmini0635-chargemagnitudequantity}
  Read a charge magnitude in a selected Physlib charge unit
\end{definition}
\begin{proof}
  This is the declared model definition of charge Readout.
\end{proof}
\begin{definition}[charge In Coulombs]
  \label{def:physics:phyx-mini-0635:phyxminiproblems-problemphyxmini0635-chargeincoulombs}
  \lean{PhyXMiniProblems.ProblemPhyXMini0635.chargeInCoulombs}
  \uses{def:physics:phyx-mini-0635:phyxminiproblems-problemphyxmini0635-chargemagnitudequantity, def:physics:phyx-mini-0635:phyxminiproblems-problemphyxmini0635-chargereadout}
  SI coulomb readout of a physical charge magnitude
\end{definition}
\begin{proof}
  This is the declared model definition of charge In Coulombs.
\end{proof}
\begin{definition}[charge In Elementary Charges]
  \label{def:physics:phyx-mini-0635:phyxminiproblems-problemphyxmini0635-chargeinelementarycharges}
  \lean{PhyXMiniProblems.ProblemPhyXMini0635.chargeInElementaryCharges}
  \uses{def:physics:phyx-mini-0635:phyxminiproblems-problemphyxmini0635-chargemagnitudequantity, def:physics:phyx-mini-0635:phyxminiproblems-problemphyxmini0635-chargereadout}
  Read a charge magnitude as a multiple of the elementary charge
\end{definition}
\begin{proof}
  This is the declared model definition of charge In Elementary Charges.
\end{proof}
\begin{definition}[speed In Meters Per Second]
  \label{def:physics:phyx-mini-0635:phyxminiproblems-problemphyxmini0635-speedinmeterspersecond}
  \lean{PhyXMiniProblems.ProblemPhyXMini0635.speedInMetersPerSecond}
  SI metres-per-second readout of a physical speed
\end{definition}
\begin{proof}
  This is the declared model definition of speed In Meters Per Second.
\end{proof}
\begin{definition}[energy In Joules]
  \label{def:physics:phyx-mini-0635:phyxminiproblems-problemphyxmini0635-energyinjoules}
  \lean{PhyXMiniProblems.ProblemPhyXMini0635.energyInJoules}
  Coherent-SI joule readout of a physical energy
\end{definition}
\begin{proof}
  This is the declared model definition of energy In Joules.
\end{proof}
\begin{definition}[physlib Elementary Charge In Coulombs]
  \label{def:physics:phyx-mini-0635:phyxminiproblems-problemphyxmini0635-physlibelementarychargeincoulombs}
  \lean{PhyXMiniProblems.ProblemPhyXMini0635.physlibElementaryChargeInCoulombs}
  Physlib's exact elementary-charge unit, expressed as a scalar number of coulombs. Charges in the setup remain dimensionful quantities
\end{definition}
\begin{proof}
  This is the declared model definition of physlib Elementary Charge In Coulombs.
\end{proof}
\begin{definition}[Figure Body]
  \label{def:physics:phyx-mini-0635:phyxminiproblems-problemphyxmini0635-figurebody}
  \lean{PhyXMiniProblems.ProblemPhyXMini0635.FigureBody}
  The two charged bodies shown in the alpha--gold diagram
\end{definition}
\begin{proof}
  This is the declared model definition of Figure Body.
\end{proof}
\begin{definition}[Interaction Stage]
  \label{def:physics:phyx-mini-0635:phyxminiproblems-problemphyxmini0635-interactionstage}
  \lean{PhyXMiniProblems.ProblemPhyXMini0635.InteractionStage}
  The two snapshots displayed in the supplied diagram
\end{definition}
\begin{proof}
  This is the declared model definition of Interaction Stage.
\end{proof}
\begin{definition}[Radial Direction]
  \label{def:physics:phyx-mini-0635:phyxminiproblems-problemphyxmini0635-radialdirection}
  \lean{PhyXMiniProblems.ProblemPhyXMini0635.RadialDirection}
  Radial direction relative to the gold nucleus
\end{definition}
\begin{proof}
  This is the declared model definition of Radial Direction.
\end{proof}
\begin{definition}[Separation Regime]
  \label{def:physics:phyx-mini-0635:phyxminiproblems-problemphyxmini0635-separationregime}
  \lean{PhyXMiniProblems.ProblemPhyXMini0635.SeparationRegime}
  Qualitative model for the initial alpha--gold separation
\end{definition}
\begin{proof}
  This is the declared model definition of Separation Regime.
\end{proof}
\begin{definition}[Interaction Model]
  \label{def:physics:phyx-mini-0635:phyxminiproblems-problemphyxmini0635-interactionmodel}
  \lean{PhyXMiniProblems.ProblemPhyXMini0635.InteractionModel}
  Interaction retained by the textbook closest-approach model
\end{definition}
\begin{proof}
  This is the declared model definition of Interaction Model.
\end{proof}
\begin{definition}[Figure Label]
  \label{def:physics:phyx-mini-0635:phyxminiproblems-problemphyxmini0635-figurelabel}
  \lean{PhyXMiniProblems.ProblemPhyXMini0635.FigureLabel}
  Literal symbolic labels visible in the primary image
\end{definition}
\begin{proof}
  This is the declared model definition of Figure Label.
\end{proof}
\begin{definition}[Alpha Gold Closest Approach Figure]
  \label{def:physics:phyx-mini-0635:phyxminiproblems-problemphyxmini0635-alphagoldclosestapproachfigure}
  \lean{PhyXMiniProblems.ProblemPhyXMini0635.AlphaGoldClosestApproachFigure}
  \uses{def:physics:phyx-mini-0635:phyxminiproblems-problemphyxmini0635-figurebody, def:physics:phyx-mini-0635:phyxminiproblems-problemphyxmini0635-interactionstage, def:physics:phyx-mini-0635:phyxminiproblems-problemphyxmini0635-figurelabel}
  Presentation facts transcribed from image 635. It contains no numerical value for r\_min
\end{definition}
\begin{proof}
  This is the declared model definition of Alpha Gold Closest Approach Figure.
\end{proof}
\begin{definition}[Alpha Gold Closest Approach Setup]
  \label{def:physics:phyx-mini-0635:phyxminiproblems-problemphyxmini0635-alphagoldclosestapproachsetup}
  \lean{PhyXMiniProblems.ProblemPhyXMini0635.AlphaGoldClosestApproachSetup}
  \uses{def:physics:phyx-mini-0635:phyxminiproblems-problemphyxmini0635-lengthquantity, def:physics:phyx-mini-0635:phyxminiproblems-problemphyxmini0635-massquantity, def:physics:phyx-mini-0635:phyxminiproblems-problemphyxmini0635-chargemagnitudequantity, def:physics:phyx-mini-0635:phyxminiproblems-problemphyxmini0635-figurebody, def:physics:phyx-mini-0635:phyxminiproblems-problemphyxmini0635-interactionstage, def:physics:phyx-mini-0635:phyxminiproblems-problemphyxmini0635-radialdirection, def:physics:phyx-mini-0635:phyxminiproblems-problemphyxmini0635-separationregime, def:physics:phyx-mini-0635:phyxminiproblems-problemphyxmini0635-interactionmodel, def:physics:phyx-mini-0635:phyxminiproblems-problemphyxmini0635-alphagoldclosestapproachfigure}
  Independent physical quantities for the two snapshots. The closest-approach distance is an unknown observable, not a definition made from an answer
\end{definition}
\begin{proof}
  This is the declared model definition of Alpha Gold Closest Approach Setup.
\end{proof}
\begin{definition}[Matches Problem Scenario And Data]
  \label{def:physics:phyx-mini-0635:phyxminiproblems-problemphyxmini0635-matchesproblemscenarioanddata}
  \lean{PhyXMiniProblems.ProblemPhyXMini0635.MatchesProblemScenarioAndData}
  \uses{def:physics:phyx-mini-0635:phyxminiproblems-problemphyxmini0635-speedinmeterspersecond, def:physics:phyx-mini-0635:phyxminiproblems-problemphyxmini0635-alphagoldclosestapproachsetup}
  The qualitative model and the incident speed stated by the problem
\end{definition}
\begin{proof}
  This is the declared model definition of Matches Problem Scenario And Data.
\end{proof}
\begin{definition}[Matches Primary Figure]
  \label{def:physics:phyx-mini-0635:phyxminiproblems-problemphyxmini0635-matchesprimaryfigure}
  \lean{PhyXMiniProblems.ProblemPhyXMini0635.MatchesPrimaryFigure}
  \uses{def:physics:phyx-mini-0635:phyxminiproblems-problemphyxmini0635-chargeinelementarycharges, def:physics:phyx-mini-0635:phyxminiproblems-problemphyxmini0635-speedinmeterspersecond, def:physics:phyx-mini-0635:phyxminiproblems-problemphyxmini0635-alphagoldclosestapproachsetup}
  Literal figure evidence and its association with the dimensionful physical charges and speeds. In particular, the image gives v\_f = 0 but does not give the value of r\_min
\end{definition}
\begin{proof}
  This is the declared model definition of Matches Primary Figure.
\end{proof}
\begin{definition}[Uses Reference Alpha Gold Constants]
  \label{def:physics:phyx-mini-0635:phyxminiproblems-problemphyxmini0635-usesreferencealphagoldconstants}
  \lean{PhyXMiniProblems.ProblemPhyXMini0635.UsesReferenceAlphaGoldConstants}
  \uses{def:physics:phyx-mini-0635:phyxminiproblems-problemphyxmini0635-massinkilograms, def:physics:phyx-mini-0635:phyxminiproblems-problemphyxmini0635-alphagoldclosestapproachsetup}
  Independent reference constants used to evaluate the multiple-choice result. The alpha-particle mass is in kilograms and Physlib's electromagnetic system method is calibrated to the free-space Coulomb constant in SI units
\end{definition}
\begin{proof}
  This is the declared model definition of Uses Reference Alpha Gold Constants.
\end{proof}
\begin{definition}[Has Physical Alpha Gold Parameters]
  \label{def:physics:phyx-mini-0635:phyxminiproblems-problemphyxmini0635-hasphysicalalphagoldparameters}
  \lean{PhyXMiniProblems.ProblemPhyXMini0635.HasPhysicalAlphaGoldParameters}
  \uses{def:physics:phyx-mini-0635:phyxminiproblems-problemphyxmini0635-lengthinmeters, def:physics:phyx-mini-0635:phyxminiproblems-problemphyxmini0635-massinkilograms, def:physics:phyx-mini-0635:phyxminiproblems-problemphyxmini0635-chargeincoulombs, def:physics:phyx-mini-0635:phyxminiproblems-problemphyxmini0635-speedinmeterspersecond, def:physics:phyx-mini-0635:phyxminiproblems-problemphyxmini0635-alphagoldclosestapproachsetup}
  Positivity conditions selecting a nondegenerate physical encounter
\end{definition}
\begin{proof}
  This is the declared model definition of Has Physical Alpha Gold Parameters.
\end{proof}
\begin{definition}[Satisfies Classical Kinetic Energy Law]
  \label{def:physics:phyx-mini-0635:phyxminiproblems-problemphyxmini0635-satisfiesclassicalkineticenergylaw}
  \lean{PhyXMiniProblems.ProblemPhyXMini0635.SatisfiesClassicalKineticEnergyLaw}
  \uses{def:physics:phyx-mini-0635:phyxminiproblems-problemphyxmini0635-massinkilograms, def:physics:phyx-mini-0635:phyxminiproblems-problemphyxmini0635-speedinmeterspersecond, def:physics:phyx-mini-0635:phyxminiproblems-problemphyxmini0635-energyinjoules, def:physics:phyx-mini-0635:phyxminiproblems-problemphyxmini0635-alphagoldclosestapproachsetup}
  Classical translational kinetic energy, K = (1/2) m v²
\end{definition}
\begin{proof}
  This is the declared model definition of Satisfies Classical Kinetic Energy Law.
\end{proof}
\begin{definition}[Satisfies Repulsive Coulomb Potential Law]
  \label{def:physics:phyx-mini-0635:phyxminiproblems-problemphyxmini0635-satisfiesrepulsivecoulombpotentiallaw}
  \lean{PhyXMiniProblems.ProblemPhyXMini0635.SatisfiesRepulsiveCoulombPotentialLaw}
  \uses{def:physics:phyx-mini-0635:phyxminiproblems-problemphyxmini0635-lengthinmeters, def:physics:phyx-mini-0635:phyxminiproblems-problemphyxmini0635-chargeincoulombs, def:physics:phyx-mini-0635:phyxminiproblems-problemphyxmini0635-energyinjoules, def:physics:phyx-mini-0635:phyxminiproblems-problemphyxmini0635-alphagoldclosestapproachsetup}
  The electrostatic energy vanishes at the effectively infinite initial separation and equals k q\_alpha q\_Au / r\_min at closest approach. This is the general Coulomb potential-energy law at the unknown separation, not the requested solution for that separation
\end{definition}
\begin{proof}
  This is the declared model definition of Satisfies Repulsive Coulomb Potential Law.
\end{proof}
\begin{definition}[Satisfies Mechanical Energy Conservation]
  \label{def:physics:phyx-mini-0635:phyxminiproblems-problemphyxmini0635-satisfiesmechanicalenergyconservation}
  \lean{PhyXMiniProblems.ProblemPhyXMini0635.SatisfiesMechanicalEnergyConservation}
  \uses{def:physics:phyx-mini-0635:phyxminiproblems-problemphyxmini0635-energyinjoules, def:physics:phyx-mini-0635:phyxminiproblems-problemphyxmini0635-alphagoldclosestapproachsetup}
  Conservation of mechanical energy between the two displayed snapshots
\end{definition}
\begin{proof}
  This is the declared model definition of Satisfies Mechanical Energy Conservation.
\end{proof}
\begin{lemma}[charge In Coulombs eq elementary Charge Multiple]
  \label{lem:physics:phyx-mini-0635:phyxminiproblems-problemphyxmini0635-chargeincoulombs-eq-elementarychargemultiple}
  \lean{PhyXMiniProblems.ProblemPhyXMini0635.chargeInCoulombs_eq_elementaryChargeMultiple}
  \uses{def:physics:phyx-mini-0635:phyxminiproblems-problemphyxmini0635-chargemagnitudequantity, def:physics:phyx-mini-0635:phyxminiproblems-problemphyxmini0635-chargeincoulombs, def:physics:phyx-mini-0635:phyxminiproblems-problemphyxmini0635-chargeinelementarycharges, def:physics:phyx-mini-0635:phyxminiproblems-problemphyxmini0635-physlibelementarychargeincoulombs}
  Unit covariance converts an elementary-charge readout to the corresponding coulomb readout. This is a unit-conversion fact, independent of the closest-approach dynamics
\end{lemma}
\begin{proof}
  The conclusion follows from the governing relations and figure readouts named in the statement of charge In Coulombs eq elementary Charge Multiple.
\end{proof}
\begin{lemma}[closest Approach Distance exact]
  \label{lem:physics:phyx-mini-0635:phyxminiproblems-problemphyxmini0635-closestapproachdistance-exact}
  \lean{PhyXMiniProblems.ProblemPhyXMini0635.closestApproachDistance_exact}
  \uses{def:physics:phyx-mini-0635:phyxminiproblems-problemphyxmini0635-lengthinmeters, def:physics:phyx-mini-0635:phyxminiproblems-problemphyxmini0635-massinkilograms, def:physics:phyx-mini-0635:phyxminiproblems-problemphyxmini0635-chargeincoulombs, def:physics:phyx-mini-0635:phyxminiproblems-problemphyxmini0635-speedinmeterspersecond, def:physics:phyx-mini-0635:phyxminiproblems-problemphyxmini0635-alphagoldclosestapproachsetup, def:physics:phyx-mini-0635:phyxminiproblems-problemphyxmini0635-matchesproblemscenarioanddata, def:physics:phyx-mini-0635:phyxminiproblems-problemphyxmini0635-matchesprimaryfigure, def:physics:phyx-mini-0635:phyxminiproblems-problemphyxmini0635-hasphysicalalphagoldparameters, def:physics:phyx-mini-0635:phyxminiproblems-problemphyxmini0635-satisfiesclassicalkineticenergylaw, def:physics:phyx-mini-0635:phyxminiproblems-problemphyxmini0635-satisfiesrepulsivecoulombpotentiallaw, def:physics:phyx-mini-0635:phyxminiproblems-problemphyxmini0635-satisfiesmechanicalenergyconservation}
  Energy conservation gives the exact closest-approach relation r\_min = 2 k q\_alpha q\_Au / (m\_alpha v\_i\textasciicircum{}2). No numerical answer choice is used in this derivation
\end{lemma}
\begin{proof}
  The conclusion follows from the governing relations and figure readouts named in the statement of closest Approach Distance exact.
\end{proof}
\begin{definition}[Answer Choice]
  \label{def:physics:phyx-mini-0635:phyxminiproblems-problemphyxmini0635-answerchoice}
  \lean{PhyXMiniProblems.ProblemPhyXMini0635.AnswerChoice}
  Labels printed beside the four multiple-choice answers
\end{definition}
\begin{proof}
  This is the declared model definition of Answer Choice.
\end{proof}
\begin{definition}[distance In Meters]
  \label{def:physics:phyx-mini-0635:phyxminiproblems-problemphyxmini0635-answerchoice-distanceinmeters}
  \lean{PhyXMiniProblems.ProblemPhyXMini0635.AnswerChoice.distanceInMeters}
  \uses{def:physics:phyx-mini-0635:phyxminiproblems-problemphyxmini0635-answerchoice}
  The closest-approach distance printed for each choice, in metres
\end{definition}
\begin{proof}
  This is the declared model definition of distance In Meters.
\end{proof}
\begin{definition}[recorded Dataset Answer]
  \label{def:physics:phyx-mini-0635:phyxminiproblems-problemphyxmini0635-recordeddatasetanswer}
  \lean{PhyXMiniProblems.ProblemPhyXMini0635.recordedDatasetAnswer}
  \uses{def:physics:phyx-mini-0635:phyxminiproblems-problemphyxmini0635-answerchoice}
  The answer label recorded in the source dataset metadata
\end{definition}
\begin{proof}
  This is the declared model definition of recorded Dataset Answer.
\end{proof}
\begin{definition}[Matches Displayed Closest Approach]
  \label{def:physics:phyx-mini-0635:phyxminiproblems-problemphyxmini0635-matchesdisplayedclosestapproach}
  \lean{PhyXMiniProblems.ProblemPhyXMini0635.MatchesDisplayedClosestApproach}
  \uses{def:physics:phyx-mini-0635:phyxminiproblems-problemphyxmini0635-lengthinmeters, def:physics:phyx-mini-0635:phyxminiproblems-problemphyxmini0635-alphagoldclosestapproachsetup, def:physics:phyx-mini-0635:phyxminiproblems-problemphyxmini0635-answerchoice}
  Agreement at the precision of a value printed to the nearest 0.1 * 10\textasciicircum{}-14 m; half a displayed unit is 5 * 10\textasciicircum{}-16 m
\end{definition}
\begin{proof}
  This is the declared model definition of Matches Displayed Closest Approach.
\end{proof}
\begin{definition}[Is Unique Matching Displayed Closest Approach]
  \label{def:physics:phyx-mini-0635:phyxminiproblems-problemphyxmini0635-isuniquematchingdisplayedclosestapproach}
  \lean{PhyXMiniProblems.ProblemPhyXMini0635.IsUniqueMatchingDisplayedClosestApproach}
  \uses{def:physics:phyx-mini-0635:phyxminiproblems-problemphyxmini0635-alphagoldclosestapproachsetup, def:physics:phyx-mini-0635:phyxminiproblems-problemphyxmini0635-answerchoice, def:physics:phyx-mini-0635:phyxminiproblems-problemphyxmini0635-matchesdisplayedclosestapproach}
  Exactly one displayed choice agrees with the physical distance
\end{definition}
\begin{proof}
  This is the declared model definition of Is Unique Matching Displayed Closest Approach.
\end{proof}
\begin{lemma}[closest Approach Distance numerical Bounds]
  \label{lem:physics:phyx-mini-0635:phyxminiproblems-problemphyxmini0635-closestapproachdistance-numericalbounds}
  \lean{PhyXMiniProblems.ProblemPhyXMini0635.closestApproachDistance_numericalBounds}
  \uses{def:physics:phyx-mini-0635:phyxminiproblems-problemphyxmini0635-lengthinmeters, def:physics:phyx-mini-0635:phyxminiproblems-problemphyxmini0635-alphagoldclosestapproachsetup, def:physics:phyx-mini-0635:phyxminiproblems-problemphyxmini0635-matchesproblemscenarioanddata, def:physics:phyx-mini-0635:phyxminiproblems-problemphyxmini0635-matchesprimaryfigure, def:physics:phyx-mini-0635:phyxminiproblems-problemphyxmini0635-usesreferencealphagoldconstants, def:physics:phyx-mini-0635:phyxminiproblems-problemphyxmini0635-hasphysicalalphagoldparameters, def:physics:phyx-mini-0635:phyxminiproblems-problemphyxmini0635-satisfiesclassicalkineticenergylaw, def:physics:phyx-mini-0635:phyxminiproblems-problemphyxmini0635-satisfiesrepulsivecoulombpotentiallaw, def:physics:phyx-mini-0635:phyxminiproblems-problemphyxmini0635-satisfiesmechanicalenergyconservation}
  With the reference alpha mass, exact elementary charge, and free-space Coulomb constant, the physical result lies strictly between 2.65 * 10\textasciicircum{}-14 m and 2.75 * 10\textasciicircum{}-14 m
\end{lemma}
\begin{proof}
  The conclusion follows from the governing relations and figure readouts named in the statement of closest Approach Distance numerical Bounds.
\end{proof}
% --- Archon declaration topology end ---
% --- Archon physics formalization source end ---
